% s-091-p091-m002; original PDF 91, 92
\Aphorism{s-091-p091-m002}{95}
Auf die besten Dinge hat die Jugend kein Recht, und auf die geringeren macht es ihr Niemand streitig. Eines unter diesen besten Dingen ist Stil; Jugend hat kein Recht auf Stil und keine Organe für Stil. Dieses Feuchte, dieser physiologische Missstand, der sich Jugend nennt, dieses ewige „Noch nicht“ und „Warte nur, balde“, mit dem sie zahlt und schuldig bleibt — unmöglich, in dieses fliessende Chaos einen Kosmos zu säen. Zum Stil gehört Reife, Kälte, Trockenheit, eine festgestellte Beziehung nach aussen und innen, eine irgendwie verbürgte Stabilität des Gleichgewichts; aber Jugend — dieser chronische Mangel an Gleichgewicht, ein Zustand, der sich jeder Formel und Definition entzieht, ein Schaffen ohne Fruchtbarkeit, ein Geniessen bis zum Verderb, ein aufgehaltenes Bergab, ein künstlich umgekehrter Rückgang — wie sollte dieses imaginäre Dasein ein Recht auf Stil haben! — Ein Volk ist stilwidrig, das sich von seiner Jugend den Ton angeben lässt. Im Alterthum war den Jünglingen unter gewissen Umständen gestattet, an den Festen der Erwachsenen theilnehmend zu lernen; wenn bei uns die Erwachsenen Feste feiern, entlehnen sie das Wie, den „Comment“, von der Jugend. In der deutschen guten Gesellschaft regirt der \Foreign{esprit} der Zwanzig- bis Fünfundzwanzigjährigen, der jungen Officiere und Studenten: Ausländer finden, dass das nicht zuviel \Foreign{esprit} sei. Die Deutschen nennen sich ein jugendliches Volk, das Volk der Zukunft; aber Jugend wird erst geniessbar, wenn sie vorbei ist.
\par

